\section{Introduction}
In the framework of the fourth industrial revolution, Digital Transformation (DX) has become a cornerstone of Vietnam’s national economic strategy. Under the "National Digital Transformation Program to 2025, with a vision to 2030" (Decision No. 749/QD-TTg), the Vietnamese government has set ambitious goals to shift from basic digitization to a comprehensive data-driven digital economy. For enterprises in Vietnam, this transition is no longer merely an option but a critical requirement to optimize operational costs and enhance global competitiveness. As a result, businesses across sectors—from banking to manufacturing—are aggressively integrating advanced technologies to restructure their traditional workflows.

Among these technologies, Robotic Process Automation (RPA) has emerged as a principal mechanism for automation in the Vietnamese market. RPA's ability to handle high-volume, repetitive tasks with speed and precision has made it an attractive entry point for enterprises from SMEs to large enterprises. By deploying software robots to mimic human interactions with digital systems, companies have successfully automated back-office functions such as payroll processing, data entry, and basic reporting.

However, as the DX journey matures, traditional RPA is reaching a significant plateau. The fundamental limitation of RPA lies in its rigid, rule-based nature: these tools are only capable of processing highly structured data (e.g., Excel or SQL databases). In the specific context of Vietnam, this limitation is magnified by several factors:

Dominance of Unstructured Data: A vast majority of business inputs in Vietnam—such as handwritten invoices, scanned delivery notes, and diverse document formats—remain unstructured.

Voice-Based Transactions: Many local SMEs still rely heavily on hand-written invoices, verbal communication and voice-based ordering, which traditional RPA cannot interpret.

Regional Linguistic Nuances: Standard automation tools often struggle with the complexity of Vietnamese dialects and non-standard shorthand in manual records.

When faced with these non-standardized and irregular inputs, traditional RPA workflows fail, necessitating human intervention and negating the benefits of end-to-end automation. To bridge this gap, Multimodal Artificial Intelligence (AI) has emerged as a transformative solution. By integrating computer vision, speech recognition, and natural language processing, Multimodal AI provides RPA with the sensory perception required to handle the complexities of the real-world Vietnamese business environment.

The primary objective of this research is to propose and implement a robust R\&D framework for an "Automated Multimodal Ordering System" specifically engineered for the Vietnamese enterprise ecosystem. Moving beyond the theoretical applications of AI, this paper focuses on the technical synergy between Computer Vision (CV), Automatic Speech Recognition (ASR) and Semantic Information Retrieval (IR) to bridge the gap between unstructured inputs and structured digital procurement records. The research aims to address the following technical challenges:

\begin{enumerate}
    \item \textbf{Architectural Design of a Modular End-to-End Pipeline}: We define a comprehensive system architecture that transforms raw Vietnamese audio signals and hand-written invoices into finalized, schema-compliant digital orders. This involves a decoupled, multi-stage process: from the selection of high-fidelity speech-to-text engines to a post-processing refinement layer and final semantic mapping. By modularizing the pipeline, the system ensures flexibility and ease of maintenance across different enterprise environments.
    \item \textbf{Strategic Model Selection and Phonetic Refinement}: A core focus of this study is the empirical evaluation and selection of Automatic Speech Recognition (ASR) models based on rigorous metrics, including Word Error Rate (WER) and Levenshtein Distance. Rather than intensive fine-tuning, we implement a Cognitive Refinement Layer using Large Language Models (LLMs) to correct common phonetic misinterpretations and dialectal inconsistencies (e.g., misspellings like "bun" for "bulb" or "oát" for "watt"). This heuristic approach significantly enhances transcription reliability in noisy, real-world commercial contexts.
    \item \textbf{Two-Tier Semantic Retrieval and Vector-Based SKU Mapping}: We explore a sophisticated information retrieval strategy that combines \textbf{LLM-based Pre-filtering} with \textbf{Vector Search mechanisms}. The LLM first cleans and extracts core entities—such as product identifiers, quantities, and units—from the refined text. These entities are then mapped against enterprise databases using high-dimensional vector embeddings, ensuring precise product identification (SKU mapping) even when the input contains non-standard terminology or regional jargon. 
\end{enumerate}